\documentclass[11pt, a4paper]{article}      
\usepackage[T2A]{fontenc}
\usepackage[utf8x]{inputenc}
\usepackage[russian]{babel}
\usepackage{csquotes}
\usepackage{amsmath}
\usepackage{amsthm}
\usepackage{enumitem}
\usepackage{amssymb}
\usepackage{listings}
\usepackage[pdftex]{graphicx}
\usepackage{caption}
\usepackage[a4paper,left=2cm,right=2cm,
top=2cm,bottom=2cm,bindingoffset=0cm]{geometry}
\usepackage[colorlinks=true, linkcolor=black, urlcolor=blue, citecolor=black]{hyperref}
\usepackage{indentfirst}
\usepackage[nottoc]{tocbibind}
\usepackage{arydshln}
\usepackage{titlesec}

\captionsetup[figure]{name=Рисунок}

\newtheorem{definition}{Определение}
\newtheorem{example}{Пример}
\newtheorem{statement}{Утверждение}
\newtheorem{theorem}{Теорема}
\newtheorem{lemm}{Лемма}

\def\intmat#1#2{\mathbb{Z}^{#1 \times #2}}
\def\ratmat#1#2{\mathbb{Q}^{#1 \times #2}}

\titlespacing*{\subsection}
{0pt}{2\baselineskip}{0.25\baselineskip}

\def\newsection#1{
    \newpage
    \section*{#1}
    \addcontentsline{toc}{section}{#1}
}
\def\newsubsection#1{
    \subsection*{#1}
    \addcontentsline{toc}{subsection}{#1}
}

\begin{document}
\setcounter{page}{1}

Пусть есть матрицы $A, B \in \intmat{n}{n}$. Множество решений уравнения $AX = XB$ над $\ratmat{n}{n}$ 
обозначим как $M_{AB} = \{X\in \ratmat{n}{n}\;|\;  AX = XB \}$. Тогда множество целочисленных решений есть $\Lambda_{AB} = M_{AB} \cap \intmat{n}{n}$.
Суть настоящей работы заключается в изучении структуры множества $\Lambda_{AB}$, а в особенности, его связь с множеством $M_{JJ}$, где 
$J$ - жорданова нормальная форма матриц $A$ и $B$.

\newsection{Свойства множества $M_{JJ}$}

Пусть есть $A, B \in \ratmat{n}{n}$ такие, что $\det(A - \lambda E) = \det(B - \lambda E) = \chi(\lambda)$ и характеристический полином $\chi(\lambda)$
раскладывается над $\mathbb{Q}$ в произведение $\chi(\lambda) = (\lambda - \lambda_1)(\lambda - \lambda_2)\cdot\ldots\cdot(\lambda - \lambda_u)$, причем среди чисел $\lambda_k$ ($k=\overline{1,u}$)
могут быть одинаковые и пусть жордановы формы этих матриц совпадают и равны $J$. Тогда $M_{JJ} = \left\{X \in \ratmat{n}{n}\; | \; JX = XJ \right\}$ и 
$M_{AB} = S_AM_{JJ}S_B^{-1}$ (Лемма 1 из \cite{Sidorov1}).

\begin{definition} \label{def:1}
    Матрицу $T$ размера $m \times n$ будем называть обобщенно-верхне-треугольной, если она имеет один из следующих видов:
    \begin{enumerate}[label={\arabic*)}]
        \item Если $m = n$, $T$ -- квадратная верхне-треугольная матрица;
        \item Если $m < n$, $T$ = $\left(\begin{matrix}
            0 & T'\\
        \end{matrix}\right)$, где $T'$ -- квадратная верхне-треугольная матрица порядка $m$;
        \item Если $m > n$, $T$ = $\left(\begin{matrix}
            T'\\
            0\\
        \end{matrix}\right)$, где $T'$ -- квадратная верхне-треугольная матрица порядка $n$;
    \end{enumerate}

    Если в каждом из трех случаев у верхне-треугольной матрицы любая наддиагональ состоит из одинаковых элементов, то такую 
    обобщенно-верхне-треугольную матрицу назовем \textbf{правильной верхне-треугольной}.

    Блочную матрицу, блоки которой представляют из себя обобщенно-верхне-треугольные (правильные верхне-треугольные) 
    матрицы будем называть \textbf{блочной обобщенно-верхне-треугольной (правильной верхне-треугольной)}.
\end{definition}

\begin{example}
    Примеры обобщенно-верхне-треугольных матриц:
    $$
    \left(\
        \begin{array}{c c c}
            0 & 1 & 2 \\
            0 & 0 & 3 \\
        \end{array}
    \right), 
    \left(\
        \begin{array}{c c }
            1 & 2 \\
            0 & 3 \\
            0 & 0 \\
        \end{array}
    \right)
    $$
    Примеры правильных верхне-треугольных матрицы:
    $$
    \left(\
        \begin{array}{c c c c}
            0 & 1 & 2 & 3 \\
            0 & 0 & 1 & 2 \\
            0 & 0 & 0 & 1 \\
        \end{array}
    \right), 
    \left(\
        \begin{array}{c c c}
            1 & 2 & 3 \\
            0 & 1 & 2 \\
            0 & 0 & 1 \\
            0 & 0 & 0 \\
        \end{array}
    \right)
    $$
    Пример блочной правильной верхне-треугольной матрицы:
    $$
    \left(\
        \begin{array}{c c | c c | c}
            1 & 2 & 0 & 0 & 1 \\
            0 & 1 & 0 & 0 & 0 \\
            \hline
            6 & 7 & 2 & 3 & 0 \\
            0 & 6 & 0 & 2 & 0 \\
            \hline
            0 & 1 & 0 & 5 & 3 \\
        \end{array}
    \right)
    $$
\end{example}

\begin{statement}[\cite{Gantmacher}, с. 199]
    Все матрицы $X \in M_{JJ}$ имеют блочно-диагональный вид $X = diag(X^{(1)}, \dots, X^{(s)})$, где $s$ -- 
    число различных собственных чисел матриц $A$ и $B$. Каждый из диагональных блоков 
    сам по себе блочная правильная верхне-треугольная матрица.
\end{statement}

\begin{example}
    Рассмотрим жорданову форму вида
    $$J =
    \left(\
        \begin{array}{c c | c c | c}
            a & 1 & 0 & 0 & 0 \\
            0 & a & 0 & 0 & 0 \\
            \hline
            0 & 0 & a & 1 & 0 \\
            0 & 0 & 0 & a & 0 \\
            \hline
            0 & 0 & 0 & 0 & a \\
        \end{array}
    \right)
    $$

    В данном случае собственное число всего одно, так что произвольная матрица $X$ из $M_{JJ}$ будет иметь вид
        $$X =
    \left(\
        \begin{array}{c c | c c | c}
            x_{11}^{(1)} & x_{11}^{(2)} & x_{12}^{(1)} & x_{12}^{(2)} & x_{13}^{(1)} \\
            0 & x_{11}^{(1)} & 0 & x_{12}^{(1)} & 0 \\
            \hline
            x_{21}^{(1)} & x_{21}^{(2)} & x_{22}^{(1)} & x_{22}^{(2)} & x_{23}^{(1)} \\
            0 & x_{21}^{(1)} & 0 & x_{22}^{(1)} & 0 \\
            \hline
            0 & x_{31}^{(1)} & 0 & x_{32}^{(1)} & x_{33}^{(1)} \\
        \end{array}
    \right)
    $$

    Обозначим $t_i$ размер $i$-го блока ($i=1,2,3$). Найдем определитель матрицы $X$ применяя теорему Лапласа. Для разложения определителя выберем столбцы с номерами
    $1, 1 + t_1, 1 + t_1 + t_2$, то есть первый столбец в каждом из блочных столбцов матрицы $X$. Не трудно проверить, что все миноры 3-го порядка лежащие на пересечении выбранных
    столбцов и строк с номерами отличными от $1, 1 + t_1, 1 + t_1 + t_2$  будут равны нулю, так как будут содержать хотя бы одну нулевую строку. Обозначим оставшийся ненулевой 
    минор $M$ и дополнительный к нему $M_1$. Тогда, в соответствии с теоремой Лапласа, $\det X = MM_1$. Выпишем эти миноры.
    $$M =
    \left|\
        \begin{array}{c c c}
            x_{11}^{(1)} & x_{12}^{(1)} & x_{13}^{(1)} \\
            x_{21}^{(1)} & x_{22}^{(1)} & x_{23}^{(1)} \\
            0 & 0 & x_{33}^{(1)} \\
        \end{array}
    \right|,
    M_1 =
    \left|\
        \begin{array}{c c}
            x_{11}^{(1)} & x_{12}^{(1)}  \\
            x_{21}^{(1)} & x_{22}^{(1)}  \\
        \end{array}
    \right|
    $$

    В свою очередь минор $M$ можно разложить по третьей строке и в результате имеем

    $$
    \det X = x_{33}^{(1)}
    \left|\
        \begin{array}{c c}
            x_{11}^{(1)} & x_{12}^{(1)}  \\
            x_{21}^{(1)} & x_{22}^{(1)}  \\
        \end{array}
    \right|^2
    $$

    То есть, мы представили определитель матрицы $X$ (в случае одного собственного числа) в виде $\det X = D_1 \cdot D_2^2$, где $D_i$ -- это некоторые миноры.
    В дальнейшем будет установлено, что порядок $i$-го минора из произведения это число блоков совпадающего порядка, а степень равна порядку этих блоков. 
\end{example}


\begin{lemm} \label{lemm:1}
    Рассмотрим случай, когда матрицы $A$ и $B$ имеют одно собственное число. Тогда всякая матрица
    $X \in M_{JJ}$ будет иметь блочную правильную верхне-треугольную структуру (Гантмахер). Обозначим блоки матрицы $X$ как
    $X_{ij}$. Размеры блока $X_{ij}$ обозначим как $t_i \times t_j$, $i, j = \overline{1, m}$ (где $m$ -- число жордановых клеток в $J$), причем
    $t_1 \geqslant t_2 \geqslant \ldots  \geqslant t_m$.
    Элемент, заполняющий главную диагональ блока $X_{ij}$ обозначим как $x^{\ast}_{ij}$.
    Пусть жорданова форма состоит из блоков $\beta$ различных порядков $s_1 > s_2 > \ldots > s_{\beta}$. Обозначим 
    через ${m_k}$ количество блоков размера $s_{k}$.

    Тогда, определитель $\det X$ имеет следующий вид:
    $$ \det X = D_{1}^{s_1} \cdot D_{2}^{s_2} \cdot \ldots \cdot D_{\beta}^{s_{\beta}}\text{,}$$
    где $D_{k}$ -- главные миноры матрицы $X$ порядка $m_{k}$ составленные из элементов $x^{\ast}_{ij}$ взятых из всех блоков матрицы 
    $X$ размера $s_{k}$.
\end{lemm}
\textit{Доказательство}.
Разложим определитель матрицы $X$ используя теорему Лапласа следующим образом: 
выберем столбцы с номерами $1, 1 + t_1, \dots, 1 + \displaystyle \sum_{k = 1}^{m}{t_k}$, и строки с аналогичными номерами. 
Полученный на пересечении выбранных строк и столбцов минор будет единственным отличным от нуля в сумме определителя. В самом деле,
Зафиксируем выбор столбцов и рассмотрим, как будут выглядеть миноры построенные на других строках. В сущности, в каждом из блоков мы 
выбрали первый столбец и можем выбирать строки. Всего возможно несколько видов блоков, в соответствии с определением \ref{def:1}. Рассмотрим
произвольный блок $X_{ij}$. Если $t_i < t_j$, то блок имеет вид $\left(\begin{matrix} 0 & X_{ij}'\\ \end{matrix}\right)$
В таком случае, очевидно, что первый столбец целиком будет состоять из нулей, а значит при любом выборе строк в минор на соответствующие позиции
будут поставлены нули. Если $t_i = t_j$, то $X_{ij}$ -- квадратная матрица верхне-треугольного вида, а значит, все элементы первого столбца, кроме, возможно, первого, 
буду равны нулю. И, наконец, если $t_i > t_j$ то блок $X_{ij} = \left(\begin{matrix} X_{ij}'\\ 0\\ \end{matrix}\right)$ и $X_{ij}'$ - квадратная верхне-треугольная матрица.
В матрице $X_{ij}'$, как мы уже выяснили, все элементы кроме первого в первом столбце обязательно равны нулю, а в оставшейся подматрице и вовсе
все элементы равны нулю. Таким образом, любой минор, стоящий на пересечении обозначенных выше столбцов и произвольных строк содержащий
хотя бы одну строку отличную от строк с номерами $1, 1 + t_1, \dots, 1 + \displaystyle \sum_{k = 1}^{m}{t_k}$ будет содержать тем самым
нулевую строку, а значит будет равен нулю.


Обозначим $M$ полученный таким образом минор и $M_1$ дополнительный к нему минор.
Тогда $\det X = M\cdot M_1$.


Покажем, что $M = D_1 \cdot D_2 \cdot \ldots \cdot D_{\beta}$.
Рассмотрим среди строк с номерами $1, 1 + t_1, \dots, 1 + \displaystyle \sum_{i = 1}^{m}{t_i}$ некоторую $s$-ую строку матрицы $X$ соответствующую блочной строке с номером $k$ ($X_{kl}, l = \overline{1, m}$).
Пусть $k_1$ наибольшее такое, что $t_{k_1} > t_k$.
Тогда, в соответствии с определением \ref{def:1} элементы, стоящие в строке $s$ на позициях с номерами $1, 1 + t_1, \dots, 1 + \displaystyle \sum_{i = 1}^{k_1 - 1}$ будут равны
нулю, так как они взяты из блоков вида $X_{kl} = \left(\begin{matrix} 0 & X_{kl}'\\ \end{matrix}\right), l=\overline{1, k_1}$ (в силу того, что $t_k < t_l$), в которых элемент, стоящий
на пересечении первой строки и первого столбца равен нулю. Начиная с позиции $p = 1 + \displaystyle \sum_{i = 1}^{k_1 + 1}{t_i}$ соответствующие блоки, содержащие элемент стоящий на пересечении $s$-ой строки и
$p$-го столбца имеют один из двух оставшихся видов, а значит тот элемент, в общем случае, отличен от нуля. Если такого номера $k_1$ нет, то такая строка не обязательно начинается с нулей. Таким образом, минор $M$ имеет ступенчатый вид:


$$M = \left|\begin{array}{c c c | c c c | c c c | c}
    x_{11}^{(1)} & \dots & x_{1n_1}^{(1)} & \dots & \dots & \dots & & & &\\
    x_{21}^{(1)} & \dots & x_{2n_1}^{(1)} & \dots & \dots & \dots & & & &\\
    \vdots & \vdots & \ddots & \vdots & \vdots & \vdots & & & &\\
    x_{n_1,1}^{(1)} & \dots & x_{n_1,n_1}^{(1)} & \dots & \dots & \dots & & &\\
    \hline
    0 & \dots & 0 & x_{n_1 + 1, n_1 + 1}^{(1)} & \dots & x_{n_1 + 1, n_2}^{(1)} & \dots & \dots & \dots & \vdots\\
    0 & \dots & 0 & x_{n_1 + 2, m_1 + 1}^{(1)} & \dots & x_{n_1 + 2, n_2}^{(1)} & \dots & \dots & \dots & \vdots\\
    \vdots & \vdots & \vdots & \vdots & \ddots & \vdots & \vdots & \vdots & \vdots & \vdots \\
    0 & \dots & 0 & x_{n_2, n_1 + 1}^{(1)} & \dots & x_{n_2, n_2}^{(1)} & \dots & \dots & \dots & \vdots\\
    \hline
    0 & \dots & 0 & 0 & \dots & 0 &  x_{n_2 + 1, n_2 + 1}^{(1)} & \dots & x_{n_2 + 1, n_3}^{(1)} & \vdots\\
    0 & \dots & 0 & 0 & \dots & 0 &  x_{n_2 + 2, n_2 + 1}^{(1)} & \dots & x_{n_2 + 2, n_3}^{(1)} & \vdots\\
    \vdots & \vdots & \vdots & \vdots & \vdots & \vdots & \vdots & \ddots & \vdots & \vdots \\
    0 & \dots & 0 & 0 & \dots & 0 &  x_{n_3, n_2 + 1}^{(1)} & \dots & x_{n_3, n_3}^{(1)} & \vdots\\
    \hline
    0 & \dots & 0 & 0 & \dots & 0 &  0 & \dots & 0 & \ast\\

\end{array}\right|$$

Заметим, что блоки стоящие на диагонали и есть $D_1, \dots, D_{\beta}$. В самом деле, исходя из того, как мы упорядочили величины $t_k$, количество нулей в начале строки
$k$-го диагонального блока будет равно $m_{k-1}$. Применяя теорему Лапласа к этому минору и получим выражение $M = D_1 \cdot D_2 \cdot \ldots \cdot D_{\beta}$.

Теперь рассмотрим вид дополнительного минора $M_1$. Этот минор получен из матрицы $X$ вычеркиванием строк и столбцов с номерами
$1, 1 + t_1, \dots, 1 + \displaystyle \sum_{i = 1}^{m}{t_i}$. То есть, в каждом блоке мы вычеркиваем первую строку и первый столбец. Полученная таким образом
матрица будет иметь такую же структуру, как и исходная матрица, однако каждый её блок станет меньше, или вовсе будет окончательно удален.
Таким образом, к минору $M_1$ можно применить тот же подход, что мы применяли к $\det X$ изначально, получив $M_1 = M^{(1)} \cdot M_2$. Минор $M^{(1)}$ будет
также как и минор $M$ раскладываться в произведение миноров $D_1, \dots, D_{\beta}$, однако, начиная с $t_{n_{\beta}} - 1$, минор с меньшим порядком $D_{\beta}$
пропадет из $M^{(t_{n_{\beta}})}$, начиная с $t_{n_{\beta - 1}} - 1$ разложение покинет минор $D_{\beta - 1}$ и так далее.
В результате, подставляя $M_l = M^{(l)} \cdot M_{l + 1}$ и раскрывая значения соответствующих миноров,
мы получим $\det X = M \cdot M_1 = M \cdot M^{(1)} \cdot M_2 = \dots = D_1^{t_{n_1}} \cdot D_2^{t_{n_2}} \cdot \ldots \cdot D_{\beta}^{t_{n_{\beta}}}$ \qed
\end{document}